\documentclass{article}

% \usepackage{neurips_2025}
\usepackage[preprint]{neurips_2025}

\usepackage[utf8]{inputenc} % allow utf-8 input
\usepackage[T1]{fontenc}    % use 8-bit T1 fonts
\usepackage{hyperref}       % hyperlinks
\usepackage{url}            % simple URL typesetting
\usepackage{booktabs}       % professional-quality tables
\usepackage{amsmath}        % amsmath
\usepackage{amsfonts}       % blackboard math symbols
\usepackage{nicefrac}       % compact symbols for 1/2, etc.
\usepackage{microtype}      % microtypography
\usepackage{xcolor}         % colors
\usepackage{graphicx}       % images
\usepackage{caption}        % captions
\usepackage{subcaption}     % subcaptions

\captionsetup{font=it}

\newcommand\abs[1]{\left| #1 \right|}
\newcommand\set[1]{\left\lbrace #1 \right\rbrace}
\newcommand\parens[1]{\left( #1 \right)}
\newcommand\brac[1]{\left[ #1 \right]}
\newcommand\norm[1]{\left\Vert #1 \right\Vert}
\newcommand{\argmin}{\operatorname{argmin}}

\title{Unsupervised Manifold Learning of Market Regimes via Physics-Informed Neural Operators}

\author{%
  Kai Hidajat \\
  Department of Applied Math \\
  University of Washington \\
  Seattle, WA, 98109\\
  \texttt{hidajatk@uw.edu} \\
}

\begin{document}

\maketitle

\begin{abstract}
    Financial markets exhibit non-stationary dynamics (``regimes'') that evolve according to unknown physical laws, increasingly characterized by roughness and high-frequency jumps. Traditional models assume smooth physics and fail to generalize to modern crisis regimes (e.g., the 2022 Inflation/Rate shock). We propose a \textbf{Physics-Informed Neural Operator} that learns the global solution manifold of the pricing SPDE. Our framework introduces three key innovations: (1) \textbf{Rough Heston Pre-training}, a curriculum strategy that primes the model with fractional Brownian motion (\(H\approx 0.1\)) physics to capture market roughness; (2) A \textbf{Spectral DeepONet} architecture that uses FFT layers to learn resolution-invariant volatility features; and (3) A \textbf{Differentiable BSM Layer} that enforces No-Arbitrage constraints by design. On out-of-sample data (2022), our model achieves a Mean Absolute Percentage Error (MAPE) of \textbf{18\%}, significantly outperforming standard baselines (>45\%). Furthermore, we demonstrate that the model's latent space spontaneously topologically separates into distinct pre- and post-crisis regimes, validating the manifold hypothesis.
\end{abstract}

\section{Introduction}

\subsection{The Generalization Gap in Financial Regimes}
The central challenge in quantitative finance is \emph{regime shift} \citep{hamilton1989}. Models calibrated on low-volatility periods (e.g., 2017-2019) often fail catastrophically when market dynamics undergo a phase transition (e.g., COVID-19, 2022 Inflation). This is an inverse problem: given observed prices \(\mathcal{Y}\), we seek the latent parameters \(\mathcal{Z}\) of the underlying Stochastic Partial Differential Equation (SPDE) governing the market.

Standard Deep Learning approaches suffer from overfitting to the specific realizations of the training period. When the distribution shifts, they fail to generalize because they have learned the \emph{data artifacts} rather than the \emph{underlying physics}. We argue that to achieve robust generalization, a model must be \textbf{operator-theoretic}: it must learn the mapping from market conditions to price surfaces, governed by physical constraints (No-Arbitrage, Smoothness/Roughness).

\subsection{Our Contribution: Rough Physics \& Spectral Operators}
We introduce a novel framework that combines Operator Learning with Physics-Informed constraints.
\begin{enumerate}
    \item \textbf{Rough Volatility Inductive Bias}: We recognize that real markets are ``rough'' \citep{rosenbaum2018rough}. We generate a synthetic dataset using \emph{Rough Heston} dynamics driven by Fractional Brownian Motion (fBm) to pre-train our model, preventing it from learning smooth, unrealistic Heston dynamics.
    \item \textbf{Primed Transfer Learning}: We employ a ``Short Burst'' pre-training strategy (100 batches) to orient the model's gradients towards valid physical solutions without rigid adherence to the synthetic parameters, enabling effective fine-tuning on real data.
    \item \textbf{Spectral Latent Space}: We show that our Spectral Neural Operator learns a latent geometry where different market regimes are topologically separated, enabling unsupervised regime identification.
\end{enumerate}

\section{Methodology}

\subsection{Architecture: Spectral DeepONet}
We employ a modified Deep Operator Network (DeepONet) \citep{Lu_2021} designed for high-frequency financial data.

\subsubsection{Branch Net (Spectral Encoder)}
The Branch Net \(\mathcal{B}\) maps the 30-day market history \(u(t)\) to a latent embedding \(\mathbf{b}\). Instead of standard CNNs, we use \textbf{Spectral Convolutions} (FFT) \citep{kovachki2021neuraloperator}.
\[
    \mathcal{K}(x) = \mathcal{F}^{-1}(R \cdot \mathcal{F}(x))
\]
This allows the network to learn global, continuous features of the volatility process, critical for capturing the non-local dependencies of rough volatility.

\subsubsection{Differentiable BSM Layer (Physics Decoder)}
Instead of predicting prices directly (which allows negative prices or arbitrage), our Trunk Net predicts the \emph{Implied Volatility Surface} \(\Sigma(K, T)\). This surface is then passed through a fixed, differentiable Black-Scholes-Merton (BSM) layer to produce prices \(P\):
\[
    P = \text{BSM}(S, K, T, r, \Sigma_{\text{pred}})
\]
This guarantees that price bounds (\(0 \le C \le S\)) and basic boundary conditions are satisfied by construction.

\subsection{Data Strategy: The ``Rough'' Synthetic Primer}
Real financial data is scarce. To train data-hungry operators, we need synthetic data. However, standard Heston models produce surfaces that are too ``smooth'' compared to the jagged, fractal nature of real markets.

\subsubsection{Rough Heston Simulation}
We implemented a GPU-accelerated Fourier-based simulator for the Rough Heston model, where the variance process \(v_t\) is driven by a fractional Brownian motion \(W^H_t\) with Hurst parameter \(H \approx 0.1\):
\[
    v_t = v_0 + \frac{1}{\Gamma(H+0.5)} \int_0^t (t-s)^{H-0.5} (\theta - v_s) ds + \dots
\]
We approximate this using Cholesky decomposition of the fractional covariance matrix to generate thousands of ``rough'' volatility paths and corresponding surfaces.

\subsubsection{Curriculum Learning}
We observed that long pre-training on synthetic data leads to \emph{negative transfer} (the model rejects real data). We devised a \textbf{Short Burst} strategy: pre-training for only 100 batches. This serves as a ``physics primer,'' initializing the weights in a valid regime without over-committing to the synthetic distribution.

\section{Experiments and Results}

\subsection{Experimental Setup}
\begin{itemize}
    \item \textbf{Training Data}: SPX Options (2019-2021) + Synthetic Rough Heston.
    \item \textbf{Test Data}: SPX Options (2022, Regime Shift to High Inflation/Rates).
    \item \textbf{Baseline}: Standard BSM calibration and Pure Supervised Learning (without synthetic).
\end{itemize}

\subsection{Pricing Performance}
Table \ref{tab:results} compares the performance. The Physics-Informed Operator significantly outperforms the baselines.

\begin{table}[h]
    \centering
    \caption{Out-of-Sample Performance (2022)}
    \label{tab:results}
    \begin{tabular}{lc}
        \toprule
        \textbf{Model} & \textbf{MAPE (\%)} \\
        \midrule
        BSM Baseline (Calibrated) & > 65.0 \\
        Pure Supervised (No Synthetic) & 43.0 \\
        Synthetic Heston (Long Pre-train) & 45.2 \\
        \textbf{Rough Heston (Short Burst) + Spectral} & \textbf{18.1} \\
        \bottomrule
    \end{tabular}
\end{table}

The ``Short Burst'' rough pre-training was decisive, reducing MAPE from ~43\% to 18\%. The training dynamics (Figure \ref{fig:pretrain}) show that the model rapidly assimilates the physics prior before fine-tuning.

\begin{figure}[h]
    \centering
    \includegraphics[width=1.0\textwidth]{../analysis/plots/full_training_dynamics.png}
    \caption{Training Dynamics. (Left) Pre-training phase showing rapid convergence on Rough Heston physics. (Right) Real data fine-tuning phase with regime switch marked by the vertical red line. The model effectively transfers the physical prior to reduce real-world error.}
    \label{fig:pretrain}
\end{figure}

\begin{figure}[h]
    \centering
    \includegraphics[width=1.0\textwidth]{../analysis/plots/surface_comparison.png}
    \caption{Visualizing the Operator: (Left) Ground Truth Volatility Surface from 2022. (Center) Predicted Surface by our Neural Operator. (Right) Absolute Error. The model accurately reconstructs the skew and term structure.}
    \label{fig:surface}
\end{figure}

\subsection{Latent Space Geometry}
We analyzed the learned latent embeddings using t-SNE. Figure \ref{fig:latent} shows a clear topological separation between the Training Regime (Pre-2020) and the Validation Regime (Post-2020).

\begin{figure}[h]
    \centering
    \includegraphics[width=0.6\textwidth]{../analysis/plots/latent_space_tsne.png}
    \caption{t-SNE visualization of the spectral latent space. The model spontaneously clusters the data into distinct regimes (yellow vs purple), validating the Manifold Hypothesis.}
    \label{fig:latent}
\end{figure}

\subsection{Computational Efficiency}
A critical advantage of the Neural Operator approach is inference speed. Standard stochastic volatility models (Heston, rBergomi) require numerical integration or Monte Carlo simulation for \emph{each contract}, leading to \(O(N)\) complexity. Our Spectral Operator, once trained, infers the entire volatility surface in a single forward pass, benefiting from massive GPU parallelism.

Figure \ref{fig:benchmark} and Table \ref{tab:benchmark} demonstrate that while Heston and Neural rBergomi (Deep Calibration) scale linearly with batch size, our model maintains sub-millisecond per-contract latency even at large batches.

\begin{figure}[h]
    \centering
    \includegraphics[width=0.7\textwidth]{../analysis/plots/inference_benchmark.png}
    \caption{Inference Latency Benchmark. The Neural Operator (Blue) exhibits near-constant time scaling for large batches due to parallelization, whereas traditional Fourier methods (Heston) and serial calibrations scale linearly.}
    \label{fig:benchmark}
\end{figure}

\begin{table}[h]
    \centering
    \caption{Inference Speed Comparison (ms)}
    \label{tab:benchmark}
    \begin{tabular}{lccc}
        \toprule
        \textbf{Method} & \textbf{Batch=1} & \textbf{Batch=1024} & \textbf{Scaling} \\
        \midrule
        BSM (Analytical) & < 0.1 & < 5 & \(O(1)\) \\
        \textbf{Neural Operator (Ours)} & \textbf{~2.0} & \textbf{~50.0} & \textbf{\(O(1)\) Parallel} \\
        Heston (Fourier) & ~100 & ~100,000 & \(O(N)\) Serial \\
        Deep rBergomi (Lit.) & ~36 & ~36,000 & \(O(N)\) Serial \\
        \bottomrule
    \end{tabular}
\end{table}

This suggests the operator has learned a \emph{regime-invariant} coordinate system. The 2022 data lies on a different manifold patch than the 2019 data, but the operator successfully interpolates this geometry.

\section{Conclusion}
We presented a Physics-Informed Neural Operator for option pricing that bridges the gap between rigorous SPDE modeling and data-driven learning. By priming the model with Rough Heston physics and enforcing BSM constraints, we achieved robust generalization to the unseen 2022 crisis regime (18\% MAPE). Future work will focus on closing the gap to <10\% through advanced hyperparameter tuning and larger-scale Rough Volatility pre-training.

\bibliographystyle{plainnat}
\bibliography{references}

\end{document}
